\section{Introduction}

The traveling salesman problem (TSP) is a classic combinatorial optimization problem that asks for the shortest possible route that visits a set of cities and returns to the origin city.
More formally, given a complete Graph \( G=(V,E) \), where \( V = \{v_{1},v_{2},\ldots ,v_n\} \) is the set of nodes representing the cities,
and \( d(v_i,v_j) \) is the distance associated with the edge between each pair of nodes \( v_i,v_j \in V \).
A solution to the TSP is an ordering of cities \( (v_{\pi(1)}, v_{\pi(2)}, \ldots  , v_{\pi(n)})\) that minimizes the total tour length:

\begin{equation}
	\left( \sum_{i=1}^{n-1} d(v_{\pi(i)}, v_{\pi(i+1)}) \right) + d(v_{\pi(n)}, v_{\pi(1)})
	\label{eq:tsp-def}
\end{equation}




Complexity: Explain why it matters (NP-hard, combinatorial explosion where brute force requires $O(n!)$ operations).


The Heuristic Objective: Explain that because exact solutions scale terribly, we trade a guarantee of absolute optimality for speed, aiming for "reasonably good" solutions
